\documentclass{article}
\usepackage[paperwidth=595.3bp,paperheight=841.9bp,margin=0bp]{geometry}
\usepackage{fontspec,eso-pic,graphicx,amsmath,booktabs,array}
\pagestyle{empty}
\setlength{\parindent}{0pt}
\setlength{\unitlength}{1bp}
\newcommand{\FontPath}{/Users/jinana/Library/Containers/com.kingsoft.wpsoffice.mac/Data/.kingsoft/office6/data/fonts/online_ext_font_cache/}
\defaultfontfeatures{Ligatures=NoCommon}
\setmainfont{Times New Roman}
\newfontfamily{\Times}{Times New Roman}
\newfontfamily{\Italic}{Times New Roman Italic}
\newfontfamily{\Bold}{Times New Roman Bold}
\newfontfamily{\BoldItalic}{Times New Roman Bold Italic}
\newfontfamily{\Sign}{KingsoftSign.ttf}[Path=/Applications/wpsoffice.app/Contents/Resources/office6/fonts/]
\newfontfamily{\Song}{HYShuSongErKW.ttf}[Path=\FontPath]
\newfontfamily{\SongHeavy}{HYShuSongErKW.ttf}[Path=\FontPath,FakeBold=2.86]
\newfontfamily{\Hei}{HYZhongHeiKW.ttf}[Path=\FontPath]
\newfontfamily{\Kai}{HYKaiTiKW.ttf}[Path=\FontPath]
\newfontfamily{\KaiHeavy}{HYKaiTiKW.ttf}[Path=\FontPath,FakeBold=2.86]
\newfontfamily{\Deng}{HYZhongDengXianKW.ttf}[Path=\FontPath]
\newcommand{\ReportPage}[1]{%
  \AddToShipoutPictureFG*{\AtPageLowerLeft{\setlength{\unitlength}{1bp}#1}}%
  \null\newpage
}
\newcommand{\Block}[4]{%
  \put(#1,\dimexpr841.9bp-#2bp\relax){\makebox(0,0)[lt]{%
  \begin{minipage}[t]{#3bp}\fontsize{10bp}{15bp}\selectfont
  \setlength{\parskip}{4bp}#4\end{minipage}}}%
}

\usepackage{xeCJK}
\usepackage{fix-cm}
\setCJKmainfont{HYShuSongErKW.ttf}[Path=\FontPath]
\begin{document}
\ReportPage{%
  \Block{90}{300}{415.3}{%
    \fontsize{9bp}{13bp}\selectfont
    记录中样品为多晶硅、非晶硅；记光强为 \(E\)，电流为 \(I\)。\par
    \vspace{-3bp}
    \includegraphics[width=415.3bp]{curves.pdf}
  }%
  \Block{90}{615}{415.3}{%
    \fontsize{9bp}{13bp}\selectfont
    \input{parameters.tex}\par
    以实测功率最大点估计 \(P_{\max}\)，\(I_{\rm SC}\) 取表3零电压点；
    \(V_{\rm OC}\) 取表2同光强读数，分别为
    \(\input{poly_voc.tex},\ \input{amorphous_voc.tex}\,\mathrm V\)。
    \[
      P_{\rm o}=VI,\quad R_m=\frac{V_m}{I_m},\quad
      FF=\frac{P_{\max}}{V_{\rm OC}I_{\rm SC}},\quad
      \eta=\frac{P_{\max}}{SE}.
    \]
    \vspace{-9bp}
    {\centering\fontsize{8.5bp}{12bp}\selectfont
    \setlength{\tabcolsep}{4bp}\input{results.tex}\par}
    \vspace{12bp}
    误差主要来自光源漂移与电池受热、探头和电池的位置及受光面积差异、
    电表内阻和接触电阻，以及峰值附近电压采样间隔较大。
    两次测量温度或光强的变化也会影响 \(FF\)。
  }%
}
\ReportPage{%
  \Block{90}{100}{415.3}{%
    暗态下两种电池均表现出整流特性：正向电流随电压加速增大，
    反向电流较小。光强增大时，短路电流近似线性增大，开路电压缓慢上升。\par
    输出电流在低电压区变化较小，接近开路时迅速下降；
    输出功率先增后减，存在最佳负载。本次相同光强和规定面积下，
    多晶硅电池的最大输出功率和转换效率均较高。
  }%
  \Block{90}{300}{415.3}{%
    1. 对伏安曲线上各点计算 \(P=VI\)，取最大值对应的
    \((V_m,I_m)\)，即 \(P_{\max}=V_mI_m\)，
    最佳匹配电阻为 \(R_m=V_m/I_m\)。
    图解上是在伏安曲线下寻找面积最大的轴平行矩形。
    应在功率峰值附近加密测点。\par
    2. 本次多晶硅电池的 \(P_{\max}=\input{poly_power.tex}\,\mathrm{mW}\)，
    非晶硅为 \(\input{amorphous_power.tex}\,\mathrm{mW}\)，
    前者约为后者的 \(\input{ratio.tex}\) 倍，主要表现为输出电流较大。
    可能原因是晶体结构、缺陷复合、吸收层厚度和载流子收集效率不同；
    电池的串并联结构、电阻损耗及对碘钨灯光谱的响应差异也会影响输出功率。
  }%
}
\end{document}
